\documentclass[12pt]{article}
\usepackage[utf8]{inputenc}
\usepackage[italian]{babel}
\author{Gabriele Guido}
\title{Emulatore CUDA}

\begin{document}
\maketitle
\clearpage
\begin{center}
    \section*{Itroduzione}
\end{center}
La rapida espansione dell'Intelligenza Artificiale ha favorito la diffusione e l'uso delle GPU 
\textit{(Graphics Processing Unit)} grazie alla loro capacità di operare in parallelo, rendendo le 
operazioni su matrici decisamente più veloci. 
La programmazione in CUDA \textit{(Compute Unified Device Architecture)} è il modo
più efficace per scrivere programmi eseguibili su GPU. L'architettura
delle schede grafiche rende la progettazione degli algoritmi meno semplice: i dati 
sono elaborati in parallelo e senza ordine quindi durante l'esecuzione potrebbero esserci delle dipendenze 
dati che, se non trattate correttamente, potrebbero portare a risultati sbagliati.
\\\\
Il progetto nasce dal bisogno di in CUDA anche senza una scheda grafica sul proprio PC. 
La libreria MyCuda contiente delle funzionalità 
necessarie alla progettazione e simulazione del calcolo parallelo. L'esecuzione 
effettiva è gestita tramite un numero limitato di thread che simulano i core di una GPU. A 
livello prestazionale non c'è alcun vantaggio, la vera utilità consiste nella verifica 
del funzionamento di algoritmi destinati al calcolo parallelo.
    
\end{document}